\subsection{The localization inside the original category}
\label{subsec:2dloc-local}

Fix $\D \in \tprl$ and a set $S$ of $1$-morphisms of $\D$.
For $(s\colon x \to y) \in S$ and $d \in \D$ let $s^{*} := \Hom_{\D}(s,d)\colon \Hom_{\D}(y,d) \to \Hom_{\D}(x,d)$, and for $f\colon d \to d'$ let $f_{*}$ be postcomposition.

\begin{defi}
\label{def:S-local}
An object $d \in \D$ is \emph{$S$-local} if every $s^{*}$, $s \in S$, admits a left adjoint $s_{!}$.
A $1$-morphism $f\colon d \to d'$ between $S$-local objects is \emph{$S$-local} if for every $s \in S$ the lax square
\[
\begin{tikzcd}[column sep=large]
\Hom_{\D}(y,d) \ar[r, "s^{*}"] \ar[d, "f_{*}"'] & \Hom_{\D}(x,d) \ar[d, "f_{*}"] \ar[dl, Rightarrow, shorten=4mm, "="'] \\
\Hom_{\D}(y,d') \ar[r, "s^{*}"'] & \Hom_{\D}(x,d')
\end{tikzcd}
\]
(given by associativity $f_{*}s^{*} \simeq s^{*}f_{*}$) is left Beck--Chevalley, i.e.\ $s_{!}f_{*} \Rightarrow f_{*}s_{!}$ is invertible.
Let $\iota\colon \Pc_{S} \hookrightarrow \D$ be the locally full sub-$(\infty,2)$-category on the $S$-local objects and $1$-morphisms.
\end{defi}

$\Pc_{S}$ is well defined: identities are $S$-local, and $S$-local $1$-morphisms are closed under composition and equivalence since mates are compatible with pasting (\cref{lem:mate-correspondence}).

\begin{thm}
\label{thm:2dloc-local}
Let $\D \in \tprl$ and let $S$ be a set of $1$-morphisms of $\D$.
\begin{enumerate}
\item $L\colon \D \to \D\radj{S}$ has a right adjoint, which is a locally full sub-$(\infty,2)$-category with image $\Pc_{S}$. Hence $\D\radj{S} \simeq \Pc_{S}$ and $L \dashv \iota$.
\item $T := \iota L$ is a $\Cat$-enriched monad on $\D$ and $\iota$ is monadic, so $\Pc_{S} \simeq \Alg_{T}(\D)$. The space of $T$-algebra structures on $d \in \D$ is empty or contractible, and non-empty if and only if $d$ is $S$-local.
\end{enumerate}
\end{thm}

\begin{proof}
(1) \emph{The underlying $\infty$-category.}
Colimits in $\tprl = \modcat{\Cat}{\prl}$ are computed in $\prl$ \cite[4.2.3.5]{ha}\todo{check the tag}, and pushouts in $\prl$ as pullbacks of right adjoints in $\widehat{\Cat}$ \cite[5.5.3.4, 5.5.3.18]{htt}.
The right adjoint of $\Free(\ell)\colon \func((\Delta^{1})\opcat,\Cat) \to \func(\Adj\opcat,\Cat)$ is restriction along $\ell\opcat$, which is $\rho$ under $(\Delta^{1})\opcat \simeq \Delta^{1}$ and $\Adj \simeq \Adj\opcat$, $l \mapsto r\opcat$, $r \mapsto l\opcat$ (in $\Adj\opcat$, $r\opcat \dashv l\opcat$). The right adjoint of $s\colon \Edge \to \D$ sends $d$ to $\Hom_{\D}(s(-),d)$, i.e.\ to $s^{*} \in \func(\Delta^{1},\Cat)$. Hence
\begin{equation}
\label{eq:pullback}
\iota_{1}\D\radj{S} \simeq \iota_{1}\D \times_{\prod_{S}\iota_{1}\func(\Delta^{1},\Cat)} \textstyle\prod_{S}\iota_{1}\func(\Adj,\Cat) \ ,
\end{equation}
with right-hand map $\prod_{S}\rho^{*}$ and projection to $\iota_{1}\D$ the right adjoint $\iota$ of $L$.
By \cref{lem:mates}(2) for $\E = \Cat$, $\rho^{*}$ is a monomorphism with image the right adjoints and the $1$-morphisms whose square is left Beck--Chevalley; for the objects $s^{*}$ these are the squares of \cref{def:S-local}. Monomorphisms are stable under products and pullbacks, so $\iota$ is a monomorphism of $\infty$-categories with image the $S$-local objects and $1$-morphisms.

\emph{The enrichment.}
By (I4), $\iota$ is enriched. As $L$ is $\Cat$-linear, $\iota$ preserves cotensors: $\mapp_{\D}(d,\iota(K \pitchfork p)) \simeq \mapp(K \otimes Ld,p) \simeq \mapp(L(K \otimes d),p) \simeq \mapp_{\D}(d,K \pitchfork \iota p)$ naturally in $d$.
For $q,p \in \D\radj{S}$ and $K \in \Cat$ we get
\[
\func\bigl(K,\Hom_{\D\radj{S}}(q,p)\bigr)^{\simeq} \simeq \mapp(q,K \pitchfork p) \subset \mapp_{\D}(\iota q,K \pitchfork \iota p) \simeq \func\bigl(K,\Hom_{\D}(\iota q,\iota p)\bigr)^{\simeq} \ ,
\]
where by the first part the left-hand side consists of the $S$-local $g\colon \iota q \to K \pitchfork \iota p$.
Now $\Hom_{\D}(s,K \pitchfork \iota p) = \func(K,\Hom_{\D}(s,\iota p))$ has left adjoint $\func(K,s_{!})$, and evaluation at $k \in K$ commutes with $s^{*}$, $s_{!}$, unit and counit. So the left mate for $g$ evaluates at $k$ to the left mate for $g_{k}\colon \iota q \to \iota p$, and $g$ is $S$-local if and only if every $g_{k}$ is, i.e.\ if and only if the corresponding functor $K \to \Hom_{\D}(\iota q,\iota p)$ factors through the full subcategory of $S$-local $1$-morphisms. By Yoneda in $K$, $\Hom_{\D\radj{S}}(q,p) \to \Hom_{\D}(\iota q,\iota p)$ is the inclusion of this full subcategory, so $\iota$ is locally fully faithful with image $\Pc_{S}$.

(2) By (I4), $T = \iota L$ is enriched with enriched $\eta$ and $\mu = \iota\varepsilon L$, and on underlying $\infty$-categories it is the monad of $L \dashv \iota$.
We check the hypotheses of Barr--Beck--Lurie \cite[4.7.3.5]{ha} for $\iota$ on underlying $\infty$-categories.
$\iota$ is conservative: if $f$ is $S$-local and invertible in $\D$, then by \cref{lem:mate-correspondence}(3) the left mate for $f^{-1}$ is inverse to that for $f$, so $f^{-1}$ is $S$-local.
$\iota$ preserves colimits of $\iota$-split simplicial objects: let $q_{\bullet}$ be a simplicial object of $\Pc_{S}$ with $\iota q_{\bullet}$ split, and write $q_{\bullet} = (\iota q_{\bullet},F_{\bullet})$ in \eqref{eq:pullback}. Colimits of split simplicial objects are preserved by every functor \cite[6.1.3.16]{htt}, so $|\iota q_{\bullet}|$ exists and $\Hom_{\D}(s,|\iota q_{\bullet}|) \simeq |\Hom_{\D}(s,\iota q_{\bullet})|$; and $|F_{\bullet}|$ exists in $\func(\Adj,\Cat)$ and is preserved by $\rho^{*}$, colimits being pointwise. So $(|\iota q_{\bullet}|,|F_{\bullet}|)$ is a colimit of $q_{\bullet}$ preserved by $\iota$.
Thus $\iota$ is monadic, $\Pc_{S} \simeq \Alg_{T}(\D)$ over $\D$, and as $\iota$ is a monomorphism the fibres of $\Alg_{T}(\D) \to \D$ are empty or contractible, non-empty exactly over the $S$-local objects.
\end{proof}

\begin{rmk}
\label{rmk:eilenberg-moore-enriched}
\Cref{thm:2dloc-local}(2) identifies $\Pc_{S}$ with $\Alg_{T}(\D)$ as $\infty$-categories. On $2$-morphisms: if $f,g\colon \iota p \to \iota q$ are $S$-local, every $\alpha\colon f \Rightarrow g$ is an algebra $2$-morphism, i.e.\ $a_{q}T(\alpha) = \alpha a_{p}$ as $2$-morphisms $a_{q}T(f) \simeq fa_{p} \Rightarrow ga_{p} \simeq a_{q}T(g)$, where $a_{p} = \iota\varepsilon_{p}$. Indeed, both sides are $2$-morphisms between $S$-local $1$-morphisms $T\iota p \to \iota q$ restricting along $\eta_{\iota p}$ to $\alpha$, and restriction along $\eta_{\iota p}$ is an equivalence $\Hom_{\Pc_{S}}(L\iota p,q) \simeq \Hom_{\D}(\iota p,\iota q)$ by the adjunction, and $\Pc_{S}$ is locally full. This is a computation in $\htwo\D$. So $T$-algebras, strong morphisms and algebra $2$-morphisms form an $(\infty,2)$-category with the same hom-categories as $\Pc_{S}$.
\todo{Compare with the Eilenberg--Moore object of \cite{riehlverity2016adjunctions,haugseng2021laxmonads}.}
\end{rmk}

\begin{rmk}
\label{rmk:2dloc-compare-bousfield}
In \cref{prop:bousfield-pushout} the local objects are those $c$ for which $s^{*}$ is an \emph{equivalence}; here $s^{*}$ is a \emph{right adjoint}.
The Beck--Chevalley condition on $1$-morphisms has no one-dimensional counterpart, since, unlike adjoints, inverses are preserved by every square.
It is what prevents $\iota$ from being fully faithful and $T$ from being idempotent: in general $L\iota p \not\simeq p$, because $\Hom_{\Pc_{S}}(L\iota p,q) \simeq \Hom_{\D}(\iota p,\iota q)$ also contains the non-local $1$-morphisms.
\end{rmk}
